\documentclass[11pt,a4paper]{article}

\usepackage[T1]{fontenc}
\usepackage[utf8]{inputenc}
\usepackage{lmodern}
\input{glyphtounicode}
\pdfgentounicode=1
\usepackage{amsmath,amssymb,amsthm}
\usepackage{mathtools}
\usepackage{hyperref}
\usepackage{doi}
\usepackage{geometry}
\usepackage{booktabs}
\usepackage{array}
\usepackage{enumitem}

\geometry{margin=2.5cm}

\hypersetup{
  hidelinks,
  pdftitle={The Non-Injective Foundations Sub-Programme: Presentation Note 5},
  pdfauthor={Jérôme Beau},
  pdfkeywords={non-injective projection, axiomatic primitive, Heisenberg group, Weil representation,
    admissibility, Born--Infeld, projection entropy, no-go theorem, emergence, foundations, presentation note}
}

\theoremstyle{definition}
\newtheorem{remark}{Remark}

\newcommand{\Heis}{\mathrm{Heis}}
\newcommand{\SU}{\mathrm{SU}}
\newcommand{\su}{\mathfrak{su}}
\newcommand{\SPI}{S_\Pi}
\newcommand{\cchi}{c_{\mathrm{BI}}}
\newcommand{\Hc}{\mathcal{H}_c}
\newcommand{\pic}{\pi_c}

\title{The Non-Injective Foundations Sub-Programme\\[0.4em]
\large Presentation Note 5}

\author{J\'{e}r\^{o}me Beau\\[0.3em]
\small Independent Researcher\\
\small \texttt{jerome.beau@cosmochrony.org}}

\date{2026\\[2pt]\small Version 1.5}

\begin{document}

  \maketitle

  \begin{abstract}
    This note presents the non-injective foundations sub-programme and types the status of
    its carrier-selection claim.  The ENI no-go theorem establishes non-injectivity as a
    structural necessity of genuine emergence, while other constituent results address
    dimensional rigidity under explicit assumptions and temporal ordering conditionally on an
    unestablished acyclicity hypothesis.  By contrast,
    the inference from admissibility axioms to
    $\Heis_3(\mathbb{Z}/q\mathbb{Z})$ and its representation data is not a theorem.
    A finite $S_3$ countermodel satisfies the algebraic properties available before the
    centrality step but has a non-central commutator.  Once the group and a non-trivial central
    character are supplied, Stone--von Neumann identifies the Heisenberg/Schr\"odinger
    representation; the associated Weil action remains distinct further symplectic data.
    Downstream results that consume these inputs are conditional on those choices.
    Conceptually, non-injectivity remains the programme's
    primitive; the selection of its physical carrier remains open.
  \end{abstract}

  \medskip
  \noindent\textbf{Keywords.}
  non-injective projection; axiomatic primitive; Heisenberg group; Weil representation;
  admissibility; Born--Infeld; projection entropy; no-go theorem; emergence; foundations;
  presentation note.

  \newpage
  \tableofcontents

  \newpage

  \section{Motivation and Central Question}
    \label{sec:motivation}

    The non-injective foundations sub-programme answers the following question.

    \medskip
    \noindent\textit{Central question.}
    Is it possible to derive $\Heis_3(\mathbb{Z}/q\mathbb{Z})$, its Schr\"odinger representation
    and associated Weil action,
    the Born--Infeld saturation constraint, and the arrow of time as consequences of a
    minimal set of structural axioms?
    More precisely:
    \textit{which structures are forced once admissible non-injective projection is taken
    as the sole primitive?}

    \medskip
    The sub-programme answers this question only in part.
    Non-injectivity is not an interpretive option or a modelling choice.
    It is a structural necessity of any genuinely emergent description (ENI).
    Within Foundation, A1+A2 give immediate information loss along non-injective transitions;
    the arrow of time follows only under the acyclicity hypothesis [H-acyc], which is not
    established.  The no-scale result supplies a separate rigidity theorem, proved under its
    stated hypotheses.
    However, A1--A3 and BI parity do not select the Heisenberg group: the finite countermodel
    of HeisenbergStructure~\cite{Beau2026HeisenbergStructure} blocks the centrality step.
    The BI capacity bound is itself the explicit axiom [A-cap].
    No background spacetime or Hilbert space is introduced at this level, but the Heisenberg
    group, Schr\"odinger carrier and associated Weil action used downstream are additional
    supplied data.

    \medskip
    This note is the definitive entry point to the corpus for a reader who wants to understand
    \emph{why} the specific structures of the other sub-programmes (spectral admissibility,
    geometry, gauge, gravity) appear, rather than merely how they are used.

  \section{Position in the Cosmochrony Programme}
    \label{sec:position}

    The non-injective foundations sub-programme is the axiomatic spine of the corpus.
    It occupies Level~1 of the three-level architecture:
    Level~1 (primitive), Level~2 (spectral fibre), Level~3 (physical observables).

    Other sub-programmes take one or more outputs or supplied realisations from this layer.
    \begin{itemize}
      \item Note~1 (Spectral Admissibility) takes $\Heis_3(\mathbb{Z}/q\mathbb{Z})$, its
      Schr\"odinger carrier $\Hc$, and the associated Weil action as separately supplied data,
      together with the BI bound $A_n \leq \cchi/\sqrt{\lambda_n}$ and the parity involution.
      \item Note~2 (Emergent Geometry) takes the supplied spectral carrier and the
      four-dimensional Carnot convergence as geometric inputs.
      On its canonical Fourier filtration the published admissibility form has the zero form as
      Mosco limit~\cite{Beau2026q5a}; the non-trivial spatial limit operator is the unestablished
      hypothesis [H-L] of Q5b~\cite{Beau2026q5b}.
      \item Note~3 (Gauge Structure) takes the phase fibre and the non-injectivity of $\Pi$ as
      the origin of gauge freedom.
      \item Note~4 (Spectral Gravity) takes $\SPI > 0$ as the structural condition that forces
      an effective gravitational sector.
    \end{itemize}

    The distinction matters: ENI is a proved reason for non-injectivity, whereas the
    Heisenberg group, its Schr\"odinger carrier and the associated Weil action are inputs to, not
    outputs proved for, the downstream calculations.

  \section{From the Preliminary Vocabulary to the Axiomatic Primitive}
    \label{sec:vocab}

    \begin{remark}[Vocabulary note: liquidation of the preliminary $\chi$ framework]
      The white paper~\cite{Cosmochrony} introduced a substrate $\chi$, a relaxation
      mechanism, and an iterated projection cascade $U_{t+1} = \Pi \circ \sigma(U_t)$ as
      the engine of dynamics.
      This vocabulary is superseded by the axiomatic formulation.

      The axiomatic formulation retains the structural content while removing the dynamical
      reading:
      the primitive is not a substrate process, but the local structure of admissible
      non-injective transitions.
      The substrate $\chi$ is static; no relaxation mechanism is postulated at the
      foundational level.
      Observable evolution is an induced ordering of admissible projections, not a
      dynamics of $\chi$ itself.
      Relaxation and iteration are not independently postulated processes in this
      formulation; where they are used as interpretations, their status must be stated at
      the point of use.

      This liquidation is final: the term ``relaxation'' should not appear in any description of
      the axiomatic framework, and $\chi$ should not be described as dynamical.
      In all notes of this series (Notes~1--9), the admissibility constraint replaces the
      relaxation vocabulary of the white paper.
    \end{remark}

  \section{Logical Chain of the Sub-Programme}
    \label{sec:chain}

    The audited logical structure is:
    \[
      \text{surjectivity, observables defined by }\Pi\text{, non-isomorphism}
      \;\Longrightarrow\;
      \underbrace{\Pi \text{ non-injective},\; \SPI > 0}_{\text{ENI, Theorem~3}}
    \]
    These are ENI's own hypotheses; ENI does not require A1--A4.  A separate route in Foundation
    obtains non-commuting admissible generators from irreducibility when the pair generates the
    acting group, the irreducibility itself being conditional on a representation-theoretic
    dictionary for $F_n$.
    The next identification does not follow:
    \[
      [X,Y]\neq e
      \;\not\Longrightarrow\;
      [X,Y]\in Z(G)
      \;\not\Longrightarrow\;
      G\simeq\Heis_3(\mathbb{Z}/q\mathbb{Z}).
    \]
    Supplying the Heisenberg group and a non-trivial central character does permit the usual
    Stone--von Neumann step to the irreducible Heisenberg/Schr\"odinger representation $\pic$ on
    $\Hc=L^2(\mathbb{Z}/q\mathbb{Z})$. It does not by itself supply the associated Weil action.

    Four status-separated statements organise the sub-programme.

    \begin{enumerate}
      \item \textit{Non-injectivity from genuine emergence (ENI~\cite{Beau2026l}).}
      Given surjective $\Pi : \Omega \to O$ with $O$ not structurally isomorphic to $\Omega$
      and no microstate-resolving structure beyond $\Pi$: $\Pi$ is necessarily non-injective
      and $\SPI > 0$.
      This is a framework-independent no-go theorem, proved from surjectivity, informational
      completeness, and non-isomorphism alone.

      \item \textit{Non-trivial commutator route (Foundation~\cite{Beau2026m}).}
      Under a representation-theoretic dictionary for the admissible fibre $F_n$, A3
      (admissibility as non-premature selection) forces its irreducibility:
      any decomposition $F_n = F^{(1)} \oplus F^{(2)}$ stable under admissible generators is
      forbidden prior to projection locking.
      On a fibre of dimension greater than one, irreducibility rules out a commuting pair of
      admissible generators when that pair generates the acting group.

      \item \textit{Carrier-selection obstruction
      (HeisenbergStructure~\cite{Beau2026HeisenbergStructure}).}
      A non-trivial commutator does not establish centrality, nilpotency, prime central order,
      or the Heisenberg presentation.  The explicit $S_3$ model satisfies the preceding
      algebraic properties while its commutator is non-central.  It therefore refutes the
      carrier-selection inference.

      \item \textit{Conditional Heisenberg representation.}
      If $\Heis_3(\mathbb{Z}/q\mathbb{Z})$ and a non-trivial central character are supplied,
      Stone--von Neumann theory identifies the corresponding irreducible Heisenberg representation
      $\pic$ on $\Hc=L^2(\mathbb{Z}/q\mathbb{Z})$. The Weil representation is a distinct associated
      symplectic action on that carrier and must be supplied or constructed through additional
      intertwiner data. The Stone--von Neumann implication survives; its antecedent is not derived.
      The parity involution $c \leftrightarrow q-c$ exists only once the central-character labels
      of that carrier are supplied, and its reading as Born--Infeld parity is an open identification;
      it does not repair the missing centrality step.
    \end{enumerate}

  \section{Constituent Papers}
    \label{sec:papers}

    \subsection{Non-injectivity as structural necessity: ENI}
      \label{sec:eni}

      \textbf{ENI}~\cite{Beau2026l} (Non-Injectivity and Strict Information Loss in Projective
      Effective Descriptions) is the foundational no-go theorem.

      \textit{Main theorem} (Theorem~3; proved).
      Any surjective $\Pi : \Omega \to O$ that is informationally complete for observables
      and such that $O$ is not structurally isomorphic to $\Omega$ must be non-injective with
      $\SPI > 0$.
      The proof uses only surjectivity, informational completeness, and the non-isomorphism
      condition; it is independent of any physical framework.

      \textit{Categorical and informational reformulations} (§8).
      Genuine emergence corresponds to a non-faithful functor (Theorem~7).
      Under the hypotheses of Theorem~3, every non-degenerate distribution $p$ on $\Omega$ has
      strictly positive projection entropy $S_\Pi(p) > 0$, which equals $H(p) - H(\Pi_* p)$ when
      both entropies are finite (Corollary~8).

      \textit{Corollary~6: Recursive non-injectivity} (proved).
      In any tower $\Omega_0 \xrightarrow{\Pi_1} \Omega_1 \xrightarrow{\Pi_2} \Omega_2$ of
      genuinely distinct emergent levels, each $\Pi_i$ is non-injective with $S_{\Pi_i} > 0$,
      and $S_{\Pi_2 \circ \Pi_1} \geq S_{\Pi_1}$.
      Non-injectivity accumulates across levels: a framework with $k$ genuinely distinct
      levels carries at least $k-1$ independent non-injective projections.

      \textit{Structural colour confinement} (§6; structural).
      The tower $\Omega_0 \to \{\text{individual sectors } c_i\}
      \to \{\text{colour-neutral pairs and triplets}\}
      \to \{\text{leptons, hadrons}\}$
      instantiates Corollary~6 at the colour stratum.
      The colour projection $\Pi_2: \{c_i\} \to \{\text{colour-neutral}\}$ is necessarily
      non-injective.
      The non-observability of free colour charges is therefore a structural consequence
      of $S_{\Pi_2} > 0$, not a dynamical accident.

    \subsection{Axiomatic foundation: Foundation~M}
      \label{sec:foundation}

      \textbf{Foundation~M}~\cite{Beau2026m} proposes a minimal axiomatic basis and develops
      several consequences.  It also states where the foundation stops short of carrier selection:
      its Theorem~5.7 is the obstruction to Heisenberg carrier selection.

      The four axioms (§2):
      \begin{itemize}
        \item A1 (Local projective admissibility): there exists a non-empty set $F_n$ of admissible
        successor directions from any observable state $O_{n-1}$.
        \item A2 (Structural non-injectivity): distinct admissible directions may lead to the same
        resolved state; $\Pi_n$ is generically non-injective.
        \item A3 (Non-premature selection): an admissible transition preserves the full multiplicity
        of open directions until projection locking; no premature selection among them is allowed.
        \item A4 (Projection locking): resolution occurs only when continuous BI saturation meets
        the discrete support of the observable image of the projection (realised as the shell
        decomposition in the spectral admissibility setting~\cite{Beau2026a26}); the transition from
        admissible to admitted is structurally discrete.
      \end{itemize}

      Results and claims, with their status:
      \begin{itemize}
        \item Immediate information loss (Proposition~4.1; proved): every admissible transition
        whose realisation map is non-injective, the generic case under A2, loses information in its
        immediate image.  Permanent non-recovery is not established (Remark~4.2).
        \item Temporal order (Theorem~4.4; conditional): along non-injective transitions, the
        observable states carry a canonical strict partial order under the acyclicity hypothesis
        [H-acyc], which is not established.  The arrow of time is therefore a conditional
        consequence, not a proved one.
        \item Structural indeterminacy of the proto-state (Proposition~5.2): from A1 and A3.
        \item Irreducibility of $F_n$ (Lemma~5.6): conditional, since A1--A3 yield it only
        together with a representation-theoretic dictionary for $F_n$.
        \item Non-commutation $[X, \sigma(X)] \neq 1$: on a fibre of dimension greater than one,
        irreducibility of the action of the group generated by $X$ and $\sigma(X)$ implies it
        (HeisenbergStructure Proposition~4.2); from A3 it is therefore conditional on the same
        dictionary.
        \item Obstruction to Heisenberg carrier selection (Theorem~5.7; proved): the
        non-commutation does not make the commutator central, so the carrier contract does
        not select $\Heis_3(\mathbb{Z}/q\mathbb{Z})$; the $S_3$ countermodel exhibits the failure.
        \item $F_n \cong \Hc$ carrying the irreducible Heisenberg representation $\pic$
        (Stone--von Neumann, detailed in HeisenbergStructure): conditional on supplying the
        Heisenberg group and central character, and on the irreducibility of $F_n$, itself
        conditional on the dictionary.  The associated Weil action is separate further
        symplectic data, not a consequence of Stone--von Neumann.
        \item Discrete character of quantum transitions (Theorem~5.12): under A3 and A4, after O22.
        \item Projective incompleteness (Corollary~4.9; proved): under A2, an admissible
        non-injective realisation map has a non-trivial fibre, hence $\ker \Pi_n \neq \{0\}$ when
        it is represented linearly.  No commutator or carrier-selection theorem is needed.
      \end{itemize}

      Three roles of $\Pi_n$ (Remark~3.1; structural):
      admissibility filter, generator of temporal order, and revelation operator.
      These roles are distinct, each irreducible to the others.  The filter role rests on the
      constraints of A1--A3 and on Corollary~4.9; the temporal role holds only under [H-acyc]
      (Theorem~4.4).

    \subsection{Finite carrier-selection obstruction: HeisenbergStructure}
      \label{sec:heisstr}
      \label{sec:hco}

      \textbf{HeisenbergStructure}~\cite{Beau2026HeisenbergStructure} audits the algebraic step from
      A1--A3 and Born--Infeld parity to $\Heis_3(\mathbb{Z}/q\mathbb{Z})$.
      It extracts the contract available before the centrality step (a finite, faithful,
      irreducible, unitary action generated by a minimal pair exchanged by an involution) and
      proves that this contract does not select a Heisenberg carrier (Theorem~4.1).
      Let $G=S_3$, let $X=(12)$ and $Y=(23)$, and let the involution be conjugation by
      $(13)$.  Then the involution exchanges $X$ and $Y$, the pair minimally generates $G$,
      and the standard two-dimensional zero-sum representation is faithful, unitary, and
      irreducible.  Nevertheless,
      \[
        [X,Y]=(132)\notin Z(S_3),
      \]
      since $Z(S_3)$ is trivial.  Thus the algebraic properties used before the centrality
      step do not select a Heisenberg group.  This is a proved obstruction to the selection
      inference, not a proposal that $S_3$ is the physical carrier.
      A separate application of Schur's lemma shows that irreducibility does not force a central
      subgroup to have prime order (Proposition~4.7).
      Stone--von Neumann remains a conditional endpoint: once the Heisenberg group and a
      non-trivial central character are supplied, it identifies $\pic$ on $\Hc$
      (Proposition~4.8); the associated Weil action is a further representation of the
      symplectic automorphism group and is not supplied by that step.
      The Robertson uncertainty bound holds on a supplied Hilbert carrier with self-adjoint
      observables (Corollary~5.1); reading it as a physical measurement principle requires an
      observable dictionary in addition to the carrier.

    \subsection{No external dimensional scale: noscale}
      \label{sec:noscale}

      \textbf{noscale}~\cite{Beau2026j} establishes a structural rigidity result:
      the admissibility structure is closed under reparameterisations of $\cchi$ but not
      under the introduction of independent dimensional parameters.

      \textit{Main theorem} (Theorem~3.1; proved).
      Any deformation of the admissibility constraints that introduces an independent
      dimensional parameter $\lambda$ (with $[\lambda] \neq [\cchi]^k$ for any $k \in \mathbb{Q}$)
      breaks at least one of:
      (i) BFS-consistent spectral scaling of $\sigma_{\mathrm{pair}}(n)$;
      (ii) structural non-injectivity (A2) and the ENI no-go;
      (iii) bounded admissibility closure (A3--A4).

      \textit{Corollary} (proved).
      Dynamical fields carrying an independent dimensional coupling cannot be consistently
      coupled to the projection $\Pi$.
      In particular, a charge $Q$ with dimensions [charge] or a cosmological constant $\Lambda$
      with dimensions [length]$^{-2}$ cannot enter the admissibility structure.
      This provides the structural backing for the exclusion of Reissner--Nordstr\"{o}m and
      Schwarzschild--de Sitter geometries in Q6b.

      The Level~2 statement --- that every emergent scale is uniquely determined as a functional
      of $\cchi$ and the cascade --- is identified as an open problem (§5 of noscale).

      \begin{table}[ht]
        \centering
        \caption{Summary of constituent papers.
        Status: P = proved; S = structural; C = conditional; supplied = model input.}
        \label{tab:papers}
        \renewcommand{\arraystretch}{1.2}
        \begin{tabular}{@{}lp{8.2cm}l@{}}
          \toprule
          Paper & Central output & Status \\
          \midrule
          ENI & Genuine emergence $\Rightarrow$ $\Pi$ non-injective $\Leftrightarrow$ $S_\Pi > 0$ & P \\
          ENI Cor.~6 & Recursive non-injectivity; colour confinement structural & P / S \\
          Foundation~M & Immediate information loss; temporal order under [H-acyc];
          carrier-selection obstruction & P / C / P \\
          HeisenbergStructure & Finite $S_3$ obstruction; Schr\"odinger representation once group and
          central character are supplied; Weil action separate & P / C / supplied \\
          noscale & No independent dimensional parameter beyond $\cchi$ & P \\
          \bottomrule
        \end{tabular}
      \end{table}

  \section{Inputs from Upstream Results}
    \label{sec:inputs}

    ENI has no upstream dependency within the corpus.
    The wider foundations synthesis is not dependency-free. It imports separately supplied
    finite-Heisenberg carrier data (group and non-trivial central character), the associated
    Weil action as further symplectic data, and the Born--Infeld candidate with its capacity axiom.

    The sole additional physical-model input beyond those supplied representation-theoretic data is
    the Born--Infeld candidate of BornInfeld~\cite{Beau2026c}: the capacity bound is stated there as the explicit
    axiom [A-cap], and the BI saturation functional is the saturation candidate
    adopted with it.
    This input is independent of A1--A4.
    The parity involution $c \leftrightarrow q-c$ is defined on the central-character labels of the
    supplied carrier; reading it as Born--Infeld parity is an open identification, since
    O18~\cite{Beau2026a22} establishes Born--Infeld parity as a covariance of the response family and leaves the fibre
    identification open.
    The Born--Infeld input neither identifies nor selects the supplied carrier from the
    foundational axioms.

  \section{Outputs to Downstream Branches}
    \label{sec:outputs}

    \begin{table}[ht]
      \centering
      \caption{Principal outputs and downstream consumers.
      Status: P = proved; S = structural; C = conditional; supplied = model input;
      Axiom = adopted axiom.}
      \label{tab:outputs}
      \renewcommand{\arraystretch}{1.3}
      \begin{tabular}{@{}p{5.8cm}p{4.2cm}l@{}}
        \toprule
        Output & Consumer & Status \\
        \midrule
        $\Pi$ non-injective; $\SPI > 0$ & Note~4 (Gravity), Q12, Q13 & P \\
        Immediate information loss along non-injective transitions & Note~2 (Geometry), Q-series & P \\
        Arrow of time (temporal order under [H-acyc]) & Note~2 (Geometry), Q-series & C \\
        $[X,Y] \neq e$ (under the representation dictionary, $\dim F_n > 1$, $\langle X,Y\rangle$
        the acting group) & Q1, Q2, HeisStr & C \\
        Heisenberg fibre group (supplied) & All Notes 1--3 & supplied \\
        $\pic$ on $\Hc=L^2(\mathbb{Z}/q\mathbb{Z})$ (Heisenberg/Schr\"odinger repr.) & Notes 1--3; Q5a--Q11 & C \\
        Associated Weil action on $\Hc$ (separate symplectic data) & Notes 1--3; O-series & supplied \\
        BI bound $A_n \leq \cchi/\sqrt{\lambda_n}$ & Note~1, O-series & Axiom \\
        Parity involution $c \leftrightarrow q-c$ (candidate fibre structure;
        conjugation identity proved, O17) & Note~1, O18, O30 & S \\
        No independent dimensional scale beyond $\cchi$ & Note~2 (noscale), Q6b & P \\
        Colour confinement as structural necessity & ENI Cor.~6, O31, O32 & S \\
        Projective incompleteness (non-injective steps, A2) & Note~3 (gauge freedom) & P \\
        \bottomrule
      \end{tabular}
    \end{table}

  \section{Status of Results}
    \label{sec:status}

    \subsection{Proved results}
      \label{sec:proved}

      \begin{enumerate}
        \item \textit{ENI Theorem~3}: genuine emergence $\Rightarrow$ $\Pi$ non-injective with
        $\SPI > 0$, proved from surjectivity, informational completeness, non-isomorphism.
        Framework-independent.

        \item \textit{ENI Corollary~6}: recursive non-injectivity; any tower of $k$ distinct
        emergent levels carries at least $k-1$ independent non-injective projections.

        \item \textit{Immediate information loss} (Foundation Proposition~4.1): a non-injective
        admissible transition, the generic case under A2, loses information in its immediate image.

        \item \textit{Non-commutation from irreducibility} (HeisenbergStructure Proposition~4.2):
        if the unitary action of $\langle X, \sigma(X)\rangle$ on $F_n$ is irreducible and
        $\dim F_n > 1$, then $X$ and $\sigma(X)$ do not commute.

        \item \textit{Finite carrier-selection obstruction}
        (HeisenbergStructure): the $S_3$ construction proves that the properties
        available before the centrality step do not imply a Heisenberg group.

        \item \textit{Discrete transitions} (Foundation Theorem~5.12): under A3 and A4, after O22.

        \item \textit{Projective incompleteness} (Foundation Corollary~4.9): under A2, a
        non-injective admissible realisation map has a non-trivial fibre.

        \item \textit{No independent dimensional parameter} (noscale Theorem~3.1):
        proved for all three cases (global multiplicative, depth-dependent, configuration-dependent
        deformations).
      \end{enumerate}

    \subsection{Structural results}
      \label{sec:structural}

      \begin{enumerate}
        \item \textit{Supplied Heisenberg carrier and Weil action}:
        the Heisenberg group, its Schr\"odinger carrier and the associated Weil action form a coherent
        and productive supplied realisation, but their selection from A1--A3 and BI parity is open.
        Results downstream of these inputs are conditional on them.

        \item \textit{Structural colour confinement} (ENI §6):
        colour confinement as a projective necessity of the emergence tower.
        The quantitative dynamics of confinement (string tension, running coupling) is not
        derived here; only the structural necessity of non-observability is established.

        \item \textit{Three roles of $\Pi_n$} (Foundation Remark~3.1):
        admissibility filter, generator of temporal order, revelation operator.
        These roles are distinct, each irreducible to the others; the temporal role holds only
        under [H-acyc].

        \item \textit{Gauge freedom as projective redundancy}:
        gauge freedom is not a separately postulated symmetry principle; it is the residual
        fibre symmetry of the non-injective projection.
        This is structurally stated in Foundation and detailed in Note~3.

        \item \textit{Level~2 scale determination} (noscale §5; open):
        that every emergent scale is uniquely determined as a functional of $\cchi$ and the
        cascade is identified as an open structural problem, beyond the Level~1 rigidity proved
        in noscale.
      \end{enumerate}

    \subsection{Conditional results}
      \label{sec:conditional}

      \begin{enumerate}
        \item \textit{Temporal order and the arrow of time} (Foundation Theorem~4.4):
        conditional on the acyclicity hypothesis [H-acyc], which is not established.

        \item \textit{Irreducibility of $F_n$} (Foundation Lemma~5.6): from A1--A3, conditional on
        a representation-theoretic dictionary for $F_n$ (HeisenbergStructure Remark~3.2).
        The non-commutation of a pair generating the acting group inherits this condition.

        \item \textit{Heisenberg representation} (HeisenbergStructure Proposition~4.8): conditional
        on supplying $\Heis_3(\mathbb{Z}/q\mathbb{Z})$ and a non-trivial central character, after
        which Stone--von Neumann theory identifies the irreducible Heisenberg representation $\pic$
        on $\Hc$.  It does not itself supply the associated Weil action.

        \item \textit{Robertson uncertainty bound} (HeisenbergStructure Corollary~5.1): conditional
        on a supplied Hilbert carrier and self-adjoint observables; its physical reading requires an
        observable dictionary.
      \end{enumerate}

    \subsection{Open hypotheses}
      \label{sec:open-hyp}

      \begin{enumerate}
        \item \textit{Carrier selection.}
        Find additional, independently motivated hypotheses that exclude the $S_3$
        countermodel and select a central extension, or state the Heisenberg group, Schr\"odinger
        carrier and Weil action as an explicit model choice.  No result in the corpus closes this bridge.

        \item \textit{Born rule.}
        Q3 derives a singlet state and correlator only on a supplied bipartite carrier and under an
        explicit diagonal-$2I$ invariance hypothesis that is not established~\cite{Beau2026q3}.
        The Born rule remains open, including in that sector; arbitrary observables,
        multipartite systems, and continuous spectra require additional structure.

        \item \textit{Full Hilbert structure from discrete fibre.}
        The passage from the discrete Schr\"odinger carrier
        $\Hc=L^2(\mathbb{Z}/q\mathbb{Z})$ to a continuous
        Hilbert space is not derived.
        Q5a proves that the published admissibility form converges to the zero form on its canonical
        Fourier filtration and, under its depth and weight hypotheses, that no common scalar
        normalisation yields a non-trivial toric differential operator~\cite{Beau2026q5a}.
        The existence of a spatial second-order limit operator is hypothesis [H-L] of
        Q5b~\cite{Beau2026q5b}, which is not established.
        The foundations sub-programme therefore treats this passage as an unclosed bridge.

        \item \textit{Level~2 scale determination (noscale)}.
        That all emergent scales (masses, coupling constants, cosmological constant) are
        uniquely determined from $\cchi$ is open.
        Only Level~1 (no independent scale can be introduced without breaking admissibility)
        is proved.
      \end{enumerate}

  \section{Open Deliverables}
    \label{sec:open}

    Three open deliverables bound the sub-programme.

    \textbf{Open deliverable 1: carrier selection.}
    Derive a physically and mathematically explicit condition that selects the
    Heisenberg group, Schr\"odinger carrier and Weil action, or retain them as a declared model
    assumption.

    \textbf{Open deliverable 2: Born rule.}
    The Born rule is not established in any sector; Q3's singlet and correlator are conditional
    on a supplied bipartite carrier and an unestablished diagonal-$2I$ invariance
    hypothesis~\cite{Beau2026q3}.
    Deriving the rule for any sector, and then for arbitrary observables and multipartite systems,
    is a primary open structural problem of this sub-programme.

    \textbf{Open deliverable 3: Level~2 scale determination.}
    Proving that every emergent scale is uniquely determined from $\cchi$ and the cascade
    (noscale Level~2) would make the framework fully predictive with a single dimensional
    free parameter.
    This requires a structural argument that no result in the corpus supplies.

  \section{Conclusion}
    \label{sec:conclusion}

    This sub-programme examines the axiomatic primitive of the Cosmochrony corpus.
    ENI establishes the structural necessity of non-injectivity under explicit hypotheses,
    while other constituent papers develop dimensional rigidity and a temporal ordering that is
    conditional on the unestablished acyclicity hypothesis [H-acyc].
    The finite $S_3$ countermodel shows that the Heisenberg group is not uniquely derived from
    A1--A3 and BI parity.  Supplying that group and a central character permits
    Stone--von Neumann to identify the Heisenberg/Schr\"odinger representation $\pic$ on $\Hc$;
    the associated Weil action remains distinct further symplectic data. Both the group input and
    the Weil action are supplied in the downstream model, so results that consume them must be read
    conditionally.

    The white-paper vocabulary of $\chi$, relaxation, and iterated projection is superseded.
    The substrate is static; the admissibility constraint replaces relaxation throughout the
    corpus.

    These results are the inputs consumed silently by every other sub-programme.
    Reading the corpus without this note is possible but the origin of each structural
    input remains opaque.
    This note makes the mixed status explicit: proved no-go results, structural
    interpretations, supplied carriers, and open identifications are distinct.  The
    conceptual lesson is not that exploration should stop, but that an identification must
    not be promoted to a theorem merely because the calculations behind it are sound.

    \newpage
    \bibliographystyle{plain}
    \bibliography{cosmochrony-bibliography}

\end{document}
